\begin{proof}[Proof of \cref{obj:thm:exact-kneser-identity}]
\begin{enumerate}
\item Let
\[
(\mathsf A_n q)(A)=\sum_{C:C\cap A=\varnothing}q(C)
\]
be the unnormalized Kneser adjacency operator on functions \(q:\Omega_{n,M}\to\mathbb R\).  For \(k=0,\ldots,M\) and \(f:\Omega_{n,M}\to\mathbb R\), the projection assumption gives \(\Pi_{n,k}f\in\mathcal H_{n,k}\), so the Kneser spectrum assumption gives
\[
\mathsf A_n(\Pi_{n,k}f)(A)
=
(-1)^k\binom{n-M-k}{M-k}\Pi_{n,k}f(A).
\]
By \cref{obj:def:kneser-covariance},
\[
\mathsf K_n(\Pi_{n,k}f)(A)
=
(-1)^k
\frac{\binom{n-M-k}{M-k}}{\binom{n-M}{M}}
\Pi_{n,k}f(A).
\]
Since \(k\le M\) and \(2M\le n\),
\[
\frac{\binom{n-M-k}{M-k}}{\binom{n-M}{M}}
=
\frac{(n-M-k)_{M-k}}{(M-k)!}\frac{M!}{(n-M)_M}
=
\frac{(M)_k}{(n-M)_k},
\]
using the falling-factorial convention \((r)_0=1\).  Hence
\[
\mathsf K_n(\Pi_{n,k}f)(A)
=
\lambda_{n,k}\Pi_{n,k}f(A).
\]
% lean: CausalSmith.Experimentation.DenseGroupPartitionProjectionPhase.exact_kneser_identity; CausalSmith.Experimentation.DenseGroupPartitionProjectionPhase.kneserChooseRatio_eq_descFactorialRatio

\item Now let \(f,g:\Omega_{n,M}\to\mathbb R\) satisfy \(\mathbb E_n f=\mathbb E_n g=0\), and put
\[
\mathcal I_+=\{j\in\{0,\ldots,M\}:j>0\}.
\]
The centered decomposition in \cref{obj:def:johnson-decomposition} gives, for every \(A\in\Omega_{n,M}\),
\[
f(A)=\sum_{j\in\mathcal I_+}\Pi_{n,j}f(A),
\qquad
g(A)=\sum_{j\in\mathcal I_+}\Pi_{n,j}g(A).
\]
For a uniform ordered disjoint pair \((A,C)\), each coordinate has the uniform \(M\)-slice marginal, and the product expectation equals the Kneser inner product:
\[
\mathbb E[f(A)]=\mathbb E_n f,\qquad
\mathbb E[g(C)]=\mathbb E_n g,
\]
and
\[
\mathbb E[f(A)g(C)]
=
\langle f,\mathsf K_n g\rangle_n .
\]
Therefore
\[
\operatorname{Cov}\{f(A),g(C)\}
=
\langle f,\mathsf K_n g\rangle_n.
\]
Substituting the positive-degree decompositions and using linearity of \(\mathsf K_n\) and of \(\langle\cdot,\cdot\rangle_n\),
\[
\begin{aligned}
\langle f,\mathsf K_n g\rangle_n
&=
\left\langle
\sum_{i\in\mathcal I_+}\Pi_{n,i}f,\,
\mathsf K_n\sum_{j\in\mathcal I_+}\Pi_{n,j}g
\right\rangle_n  \\
&=
\sum_{i\in\mathcal I_+}\sum_{j\in\mathcal I_+}
\lambda_{n,j}\,
\langle \Pi_{n,i}f,\Pi_{n,j}g\rangle_n  \\
&=
\sum_{j\in\mathcal I_+}
\lambda_{n,j}\,
\langle \Pi_{n,j}f,\Pi_{n,j}g\rangle_n ,
\end{aligned}
\]
because distinct Johnson degrees are orthogonal.  This proves the covariance identity.
% lean: CausalSmith.Experimentation.DenseGroupPartitionProjectionPhase.exact_kneser_identity; CausalSmith.Experimentation.DenseGroupPartitionProjectionPhase.orderedDisjointPair_first_E_eq_slice; CausalSmith.Experimentation.DenseGroupPartitionProjectionPhase.orderedDisjointPair_second_E_eq_slice; CausalSmith.Experimentation.DenseGroupPartitionProjectionPhase.orderedDisjointPair_E_eq

\item For each arm \(z\in\{0,1\}\), define the centered arm table
\[
h_{z,n}^{\circ}(A)=h_{z,n}(A)-\mathbb E_n h_{z,n}.
\]
Then \(\mathbb E_n h_{z,n}^{\circ}=0\).  By \cref{obj:def:kneser-covariance} and the preceding ordered-pair identity,
\[
C_{ab,n}
=
\operatorname{Cov}\{h_{a,n}^{\circ}(A),h_{b,n}^{\circ}(C)\}
\]
for \(a,b\in\{0,1\}\) under a uniform ordered disjoint pair \((A,C)\).  Applying the covariance identity from the previous step to \((h_{1,n}^{\circ},h_{1,n}^{\circ})\), \((h_{0,n}^{\circ},h_{0,n}^{\circ})\), and \((h_{1,n}^{\circ},h_{0,n}^{\circ})\), and using
\[
C_{\tau\tau,n}=C_{11,n}+C_{00,n}-2C_{10,n},
\]
gives
\[
C_{\tau\tau,n}
=
\sum_{j\in\mathcal I_+}\lambda_{n,j}
\left[
\langle \Pi_{n,j}h_{1,n}^{\circ},\Pi_{n,j}h_{1,n}^{\circ}\rangle_n
+
\langle \Pi_{n,j}h_{0,n}^{\circ},\Pi_{n,j}h_{0,n}^{\circ}\rangle_n
-
2\langle \Pi_{n,j}h_{1,n}^{\circ},\Pi_{n,j}h_{0,n}^{\circ}\rangle_n
\right].
\]
% lean: CausalSmith.Experimentation.DenseGroupPartitionProjectionPhase.exact_kneser_identity

\item It remains to identify the summand with the squared norm of the projected arm contrast.  For every positive degree \(j\in\mathcal I_+\) and constant \(r\in\mathbb R\),
\[
\Pi_{n,j}(r)=0.
\]
Indeed, the centered decomposition applied to the constant function \(r\) gives
\[
\sum_{\ell\in\mathcal I_+}\Pi_{n,\ell}(r)=0.
\]
Taking the slice inner product with \(\Pi_{n,j}(r)\) and using degree orthogonality leaves
\[
\langle \Pi_{n,j}(r),\Pi_{n,j}(r)\rangle_n=0,
\]
hence \(\Pi_{n,j}(r)=0\) by definiteness of the slice inner product.  Therefore, by linearity of \(\Pi_{n,j}\),
\[
\Pi_{n,j}h_{z,n}^{\circ}
=
\Pi_{n,j}\bigl(h_{z,n}-\mathbb E_n h_{z,n}\bigr)
=
\Pi_{n,j}h_{z,n}.
\]
Again by linearity,
\[
\Pi_{n,j}(h_{1,n}-h_{0,n})
=
\Pi_{n,j}h_{1,n}-\Pi_{n,j}h_{0,n}.
\]
Thus
\[
\begin{aligned}
\left\|\Pi_{n,j}(h_{1,n}-h_{0,n})\right\|_n^2
&=
\left\langle
\Pi_{n,j}h_{1,n}-\Pi_{n,j}h_{0,n},\,
\Pi_{n,j}h_{1,n}-\Pi_{n,j}h_{0,n}
\right\rangle_n  \\
&=
\langle \Pi_{n,j}h_{1,n},\Pi_{n,j}h_{1,n}\rangle_n
+
\langle \Pi_{n,j}h_{0,n},\Pi_{n,j}h_{0,n}\rangle_n
-
2\langle \Pi_{n,j}h_{1,n},\Pi_{n,j}h_{0,n}\rangle_n .
\end{aligned}
\]
Substitution into the preceding display yields
\[
C_{\tau\tau,n}
=
\sum_{j=1}^{M}\lambda_{n,j}
\bigl\|\Pi_{n,j}(h_{1,n}-h_{0,n})\bigr\|_n^2,
\]
which is the asserted contrast identity.
% lean: CausalSmith.Experimentation.DenseGroupPartitionProjectionPhase.johnson_proj_centered_eq; CausalSmith.Experimentation.DenseGroupPartitionProjectionPhase.sliceInner_sub_self; CausalSmith.Experimentation.DenseGroupPartitionProjectionPhase.exact_kneser_identity
\end{enumerate}
\end{proof}

\begin{proof}[Proof of \cref{obj:thm:exact-pame-variance}]
By the group-size hypotheses, \(G_{1n}>0\), \(G_{1n}<G_n\), \(G_{0n}=G_n-G_{1n}>0\), and \(2\le G_n\). The feasibility condition and \(2\le G_n\) also give \(2M\le n\). Write the treated-group fraction as
\[
p_n=\frac{G_{1n}}{G_n},
\qquad
q_n=1-p_n=\frac{G_{0n}}{G_n}.
\]
For a realized ordered partition
\(\boldsymbol\alpha=(\alpha_1,\ldots,\alpha_{G_n})\) of pairwise-disjoint elements of
\(\Omega_{n,M}\), define
\[
x_{zg}(\boldsymbol\alpha)=h_{z,n}(\alpha_g),\qquad
d_g(\boldsymbol\alpha)=x_{1g}(\boldsymbol\alpha)-x_{0g}(\boldsymbol\alpha),
\]
and, for any vector \(\mathbf u=(u_1,\ldots,u_{G_n})\in\mathbb R^{G_n}\),
\[
\bar u=\frac1{G_n}\sum_{g=1}^{G_n}u_g,
\qquad
\mathsf S(\mathbf u)=\frac1{G_n-1}\sum_{g=1}^{G_n}(u_g-\bar u)^2 .
\]
Set
\[
\mathsf S_z(\boldsymbol\alpha)=\mathsf S\bigl((x_{zg}(\boldsymbol\alpha))_{g=1}^{G_n}\bigr),
\qquad
\mathsf S_\tau(\boldsymbol\alpha)=\mathsf S\bigl((d_g(\boldsymbol\alpha))_{g=1}^{G_n}\bigr).
\]
These quantities are well defined because \(G_n\ge2\). % lean: exact_pame_variance

Let
\[
\Delta_n(A)=h_{1,n}(A)-h_{0,n}(A),
\qquad
\Delta_n^\circ(A)=\Delta_n(A)-\mathbb E_n\Delta_n(A).
\]
Since \(h_{z,n}^\circ(A)=h_{z,n}(A)-\mathbb E_n h_{z,n}(A)\), we have
\[
\Delta_n^\circ=h_{1,n}^\circ-h_{0,n}^\circ .
\]
Using the bilinearity of the uniform-slice inner product in \cref{obj:def:kneser-covariance} and the symmetry of the uniform ordered-disjoint-pair law under swapping the two groups,
\[
\begin{aligned}
\left\langle \Delta_n^\circ,\mathsf K_n\Delta_n^\circ\right\rangle_n
&=
C_{11,n}+C_{00,n}-C_{10,n}-C_{01,n}  \\
&=
C_{11,n}+C_{00,n}-2C_{10,n}
=
C_{\tau\tau,n}.
\end{aligned}
\]
% lean: disjointCov_sub_self

We next record the two partition-moment identities used below. For any
\(\phi:\Omega_{n,M}\to\mathbb R\), let
\[
\mu_\phi=\mathbb E_n\phi(A),\qquad
\eta_\phi=\mathbb E_n\{\phi(A)^2\},
\]
and let
\[
\kappa_\phi
=
\mathbb E_A\mathbb E_{C\mid C\cap A=\varnothing}\{\phi(A)\phi(C)\}.
\]
Then
\[
\operatorname{Var}_n(\phi)=\eta_\phi-\mu_\phi^2,
\qquad
\left\langle \phi-\mu_\phi,\mathsf K_n(\phi-\mu_\phi)\right\rangle_n
=\kappa_\phi-\mu_\phi^2 .
\]
Each coordinate of \(\boldsymbol\alpha\) is uniform on \(\Omega_{n,M}\), and each distinct ordered coordinate pair is a uniform ordered-disjoint pair. Therefore
\[
\begin{aligned}
\mathbb E_{\boldsymbol\alpha}\mathsf S\bigl((\phi(\alpha_g))_{g=1}^{G_n}\bigr)
&=
\frac{1}{G_n-1}
\left\{
G_n\eta_\phi
-\frac1{G_n}\bigl(G_n\eta_\phi+G_n(G_n-1)\kappa_\phi\bigr)
\right\} \\
&=\eta_\phi-\kappa_\phi \\
&=
\operatorname{Var}_n(\phi)
-
\left\langle \phi-\mu_\phi,\mathsf K_n(\phi-\mu_\phi)\right\rangle_n .
\end{aligned}
\]
Applying this with \(\phi=h_{1,n}\), \(\phi=h_{0,n}\), and \(\phi=\Delta_n\) gives
\[
\mathbb E_{\boldsymbol\alpha}\mathsf S_1=V_{1,n}-C_{11,n},\qquad
\mathbb E_{\boldsymbol\alpha}\mathsf S_0=V_{0,n}-C_{00,n},
\]
and
\[
\mathbb E_{\boldsymbol\alpha}\mathsf S_\tau
=
\operatorname{Var}_n(\Delta_n)-C_{\tau\tau,n}.
\]
% lean: partition_expected_sample_variance

The same exchangeability calculation gives the variance of the partition average. Indeed,
\[
\operatorname{Var}_{\boldsymbol\alpha}
\left\{
\frac1{G_n}\sum_{g=1}^{G_n}\phi(\alpha_g)
\right\}
=
\frac{1}{G_n^2}
\left\{
G_n\operatorname{Var}_n(\phi)
+
G_n(G_n-1)
\left\langle \phi-\mu_\phi,\mathsf K_n(\phi-\mu_\phi)\right\rangle_n
\right\}.
\]
Thus, for \(\phi=\Delta_n\),
\[
\operatorname{Var}_{\boldsymbol\alpha}
\left\{
\frac1{G_n}\sum_{g=1}^{G_n}d_g(\boldsymbol\alpha)
\right\}
=
\frac{\operatorname{Var}_n(\Delta_n)+(G_n-1)C_{\tau\tau,n}}{G_n}.
\]
Also, since each coordinate \(\alpha_g\) is uniform on \(\Omega_{n,M}\),
\[
\mathbb E_{\boldsymbol\alpha}
\left\{
\frac1{G_n}\sum_{g=1}^{G_n}d_g(\boldsymbol\alpha)
\right\}
=
\mathbb E_n h_{1,n}-\mathbb E_n h_{0,n}.
\]
% lean: partition_mean_variance

Condition on \(\boldsymbol\alpha\). Under the balanced assignment in
\cref{obj:def:random-partition-design},
\[
\mathbb E(Z_g\mid\boldsymbol\alpha)=p_n,\qquad
\mathbb E(1-Z_g\mid\boldsymbol\alpha)=q_n,
\]
and the usual complete-randomization second moments for distinct \(g\ne \ell\) are
\[
\mathbb E(Z_gZ_\ell\mid\boldsymbol\alpha)
=\frac{G_{1n}(G_{1n}-1)}{G_n(G_n-1)},
\]
\[
\mathbb E\{(1-Z_g)(1-Z_\ell)\mid\boldsymbol\alpha\}
=\frac{G_{0n}(G_{0n}-1)}{G_n(G_n-1)},
\]
\[
\mathbb E\{Z_g(1-Z_\ell)\mid\boldsymbol\alpha\}
=\frac{G_{1n}G_{0n}}{G_n(G_n-1)}.
\]
Expanding the difference-in-means estimator from \cref{obj:def:pame-estimator} with these moments gives
\[
\mathbb E(\widehat\tau_n\mid\boldsymbol\alpha)
=
\frac1{G_n}\sum_{g=1}^{G_n}d_g(\boldsymbol\alpha),
\]
and
\[
\operatorname{Var}(\widehat\tau_n\mid\boldsymbol\alpha)
=
\frac{\mathsf S_1(\boldsymbol\alpha)}{G_{1n}}
+
\frac{\mathsf S_0(\boldsymbol\alpha)}{G_{0n}}
-
\frac{\mathsf S_\tau(\boldsymbol\alpha)}{G_n}.
\]
The same complete-randomization calculation for the denominator-\(G_{zn}-1\) within-arm sample variances in \cref{obj:def:cr2-statistic} yields
\[
\mathbb E(\widehat V_{\mathrm{CR2},n}\mid\boldsymbol\alpha)
=
\frac{\mathsf S_1(\boldsymbol\alpha)}{G_{1n}}
+
\frac{\mathsf S_0(\boldsymbol\alpha)}{G_{0n}} .
\]
% lean: conditional_Var_pameHat

Taking expectations in the conditional mean identity and using the partition-average identity above,
\[
\mathbb E[\widehat\tau_n]
=
\mathbb E_n h_{1,n}-\mathbb E_n h_{0,n}
=
\tau_n,
\]
where the last equality is \cref{obj:def:pame-estimator}. % lean: conditional_E_pameHat

By the finite tower variance identity for the two-stage design,
\[
\sigma_n^2
=
\mathbb E_{\boldsymbol\alpha}\{\operatorname{Var}(\widehat\tau_n\mid\boldsymbol\alpha)\}
+
\operatorname{Var}_{\boldsymbol\alpha}\{\mathbb E(\widehat\tau_n\mid\boldsymbol\alpha)\}.
\]
Substituting the identities already proved gives
\[
\begin{aligned}
\sigma_n^2
&=
\frac{V_{1,n}-C_{11,n}}{G_{1n}}
+
\frac{V_{0,n}-C_{00,n}}{G_{0n}}
-
\frac{\operatorname{Var}_n(\Delta_n)-C_{\tau\tau,n}}{G_n}
+
\frac{\operatorname{Var}_n(\Delta_n)+(G_n-1)C_{\tau\tau,n}}{G_n} \\
&=
\frac{V_{1,n}-C_{11,n}}{G_{1n}}
+
\frac{V_{0,n}-C_{00,n}}{G_{0n}}
+
C_{\tau\tau,n}.
\end{aligned}
\]
Since
\[
\frac{R_n}{G_n}
=
\frac1{G_n}
\left\{
\frac{V_{1,n}}{G_{1n}/G_n}
+
\frac{V_{0,n}}{G_{0n}/G_n}
\right\}
=
\frac{V_{1,n}}{G_{1n}}+\frac{V_{0,n}}{G_{0n}},
\]
this is exactly
\[
\sigma_n^2
=
\frac{R_n}{G_n}
+
C_{\tau\tau,n}
-
\frac{C_{11,n}}{G_{1n}}
-
\frac{C_{00,n}}{G_{0n}} .
\]
% lean: finiteDesign_Var_compoundCore_tower

Taking expectations in the conditional CR2 identity gives
\[
\begin{aligned}
\mathbb E[\widehat V_{\mathrm{CR2},n}]
&=
\frac{\mathbb E_{\boldsymbol\alpha}\mathsf S_1(\boldsymbol\alpha)}{G_{1n}}
+
\frac{\mathbb E_{\boldsymbol\alpha}\mathsf S_0(\boldsymbol\alpha)}{G_{0n}} \\
&=
\frac{V_{1,n}-C_{11,n}}{G_{1n}}
+
\frac{V_{0,n}-C_{00,n}}{G_{0n}} .
\end{aligned}
\]
% lean: conditional_E_cr2Var

Subtracting the variance identity from the CR2 expectation identity, the common terms cancel and therefore
\[
\mathbb E[\widehat V_{\mathrm{CR2},n}]-\sigma_n^2
=
-C_{\tau\tau,n}.
\]
% lean: exact_pame_variance
\end{proof}

\begin{proof}[Proof of \cref{obj:thm:dense-projection-limit}]
Set
\[
\Gamma:=64M^2(M+1)B^2,
\qquad
C:=\Gamma+1 .
\]
Then \(C>0\).  Fix a feasible triangular array satisfying the four displayed assumptions, fix the rowwise Johnson projections, and bring into scope the two rowwise hypotheses that these projections give the centered orthogonal Johnson decomposition and that the corresponding Kneser spectrum formula holds.  Index rows by \(r\), and write \(n_r\) for the row population size, \(G_r\) for the group count, \(G_{1r}\) and \(G_{0r}=G_r-G_{1r}\) for the arm counts, \(N_r=MG_r\), and \(p_r=G_{1r}/G_r\).  Since \(2\le M\), \(M>0\).  The feasibility condition gives
\[
G_r\le MG_r=N_r\le n_r .
\]
Together with \cref{obj:ass:group-count-growth}, this implies \(n_r\to\infty\).  Hence
\[
\frac1{n_r}\to0,\qquad
\frac{M}{n_r}\to0,\qquad
1-\frac{M}{n_r}\to1 .
\]
% lean: CausalSmith.Experimentation.DenseGroupPartitionProjectionPhase.dense_projection_limit

Define
\[
\alpha_r:=\frac{MG_r}{n_r-M}.
\]
Because \(G_{1r}>0\) and \(G_{1r}<G_r\), every row has \(G_r\ge2\), and therefore \(2M\le MG_r\le n_r\), so the denominator is positive.  Also
\[
\alpha_r
=
\frac{MG_r/n_r}{1-M/n_r}
=
\frac{N_r/n_r}{1-M/n_r}.
\]
By \cref{obj:ass:sampling-fraction} and the preceding display, \(\alpha_r\to\rho\).
% lean: CausalSmith.Experimentation.DenseGroupPartitionProjectionPhase.dense_projection_limit

For \(S\in\Omega_{n_r,M}\), \cref{obj:def:group-table} and \cref{obj:ass:bounded-potential-outcomes} give, for each arm \(z\),
\[
|h_{z,r}(S)|
=
\left|
\frac1M\sum_{i\in S}Y_{i,r}(z,S)
\right|
\le
\frac1M\sum_{i\in S}|Y_{i,r}(z,S)|
\le B .
\]
Put
\[
d_r(S):=h_{1,r}(S)-h_{0,r}(S).
\]
Then
\[
|d_r(S)|\le |h_{1,r}(S)|+|h_{0,r}(S)|\le 2B .
\]
More generally, if \(|u(S)|\le L\) on \(\Omega_{n_r,M}\), then every Johnson projection satisfies
\[
\|\Pi_{r,k}u\|_r^2\le L^2 .
\]
Indeed, writing \(v=\Pi_{r,k}u\), orthogonal projection gives
\[
\|u\|_r^2=\|u-v\|_r^2+\|v\|_r^2,
\]
so \(\|v\|_r^2\le \|u\|_r^2=\mathbb E_r[u(S)^2]\le L^2\).  Applying this with \(u=d_r\), \(L=2B\), and \(k=1\) yields
\[
0\le E_{\tau,1,r}
=
\|\Pi_{r,1}d_r\|_r^2
\le (2B)^2 .
\]
Consequently
\[
(\rho-\alpha_r)E_{\tau,1,r}\to0 .
\]
% lean: CausalSmith.Experimentation.DenseGroupPartitionProjectionPhase.johnsonProjection_energy_le_of_abs_le

Let
\[
\mathcal I_{\ge2}:=\{k\in\{0,\ldots,M\}:2\le k\}
\]
and
\[
H_r:=
G_r\sum_{k\in\mathcal I_{\ge2}}\lambda_{r,k}\|\Pi_{r,k}d_r\|_r^2 ,
\]
where \(\lambda_{r,k}\) is the normalized Kneser eigenvalue from \cref{obj:thm:exact-kneser-identity}.  We claim that, for every row,
\[
|H_r|\le \frac{\Gamma}{n_r}.
\]
First suppose \(4M\le n_r\).  Put \(\eta_r:=2M/n_r\), so \(0\le\eta_r\le1\).  For any function \(u\) on \(\Omega_{n_r,M}\), any finite degree set \(\mathcal I\subseteq\{0,\ldots,M\}\), any integer threshold \(q\) such that \(q\le k\) for all \(k\in\mathcal I\), and any envelope \(L\) satisfying \(|u(S)|\le L\) for all \(S\), the spectral-sum estimate is
\[
\left|
\sum_{k\in\mathcal I}\lambda_{r,k}\|\Pi_{r,k}u\|_r^2
\right|
\le
|\mathcal I|\eta_r^qL^2 .
\]
Indeed, from
\[
\lambda_{r,k}=(-1)^k\frac{(M)_k}{(n_r-M)_k},
\]
each numerator factor is at most \(M\), while each denominator factor is at least \(n_r/2\); hence
\[
|\lambda_{r,k}|\le\eta_r^k\le\eta_r^q
\qquad (k\in\mathcal I).
\]
The projection-energy bound proved above gives
\[
\|\Pi_{r,k}u\|_r^2\le L^2,
\]
and the triangle inequality therefore gives the displayed spectral-sum estimate.
% lean: CausalSmith.Experimentation.DenseGroupPartitionProjectionPhase.johnsonSpectralSum_abs_le

Applying this estimate with \(\mathcal I=\mathcal I_{\ge2}\), \(q=2\), \(u=d_r\), and \(L=2B\) gives
\[
\left|
\sum_{k\in\mathcal I_{\ge2}}\lambda_{r,k}\|\Pi_{r,k}d_r\|_r^2
\right|
\le
(M+1)\left(\frac{2M}{n_r}\right)^2(2B)^2 .
\]
Using \(G_r\le n_r\),
\[
|H_r|
\le
n_r(M+1)\left(\frac{2M}{n_r}\right)^2(2B)^2
=
\frac{16M^2(M+1)B^2}{n_r}
\le
\frac{\Gamma}{n_r}.
\]
Now suppose \(4M>n_r\).  Since \(2M\le n_r\), the same eigenvalue formula gives
\[
|\lambda_{r,k}|=\frac{(M)_k}{(n_r-M)_k}\le1
\qquad (0\le k\le M),
\]
because \(M\le n_r-M\).  Together with \(\|\Pi_{r,k}d_r\|_r^2\le(2B)^2\), this yields
\[
\left|
\sum_{k\in\mathcal I_{\ge2}}\lambda_{r,k}\|\Pi_{r,k}d_r\|_r^2
\right|
\le
(M+1)(2B)^2 .
\]
Using \(G_r\le n_r\) and \(n_r\le4M\),
\[
|H_r|
\le
n_r(M+1)(2B)^2
=4n_r(M+1)B^2
\le
\frac{64M^2(M+1)B^2}{n_r}
=
\frac{\Gamma}{n_r}.
\]
Since \(n_r\to\infty\), \(H_r\to0\).
% lean: CausalSmith.Experimentation.DenseGroupPartitionProjectionPhase.higherDegreeContribution_uniform_bound

By \cref{obj:thm:exact-kneser-identity}, applied to \(d_r\),
\[
C_{\tau\tau,r}
=
\sum_{k=1}^{M}\lambda_{r,k}\|\Pi_{r,k}d_r\|_r^2 .
\]
The degree-one eigenvalue is
\[
\lambda_{r,1}=-\frac{M}{n_r-M}.
\]
Splitting off degree one and multiplying by \(G_r\) gives
\[
G_rC_{\tau\tau,r}
=
-\frac{MG_r}{n_r-M}E_{\tau,1,r}
+
G_r\sum_{k=2}^{M}\lambda_{r,k}\|\Pi_{r,k}d_r\|_r^2
=
-\alpha_r E_{\tau,1,r}+H_r .
\]
Therefore
\[
G_rC_{\tau\tau,r}+\rho E_{\tau,1,r}
=
(\rho-\alpha_r)E_{\tau,1,r}+H_r
\to0 .
\]
This proves the second displayed limit.
% lean: CausalSmith.Experimentation.DenseGroupPartitionProjectionPhase.scaledContrastCrossCov_decomposition

Next fix an arm \(z\).  Let
\[
h_{z,r}^{\circ}(S):=h_{z,r}(S)-\mathbb E_r h_{z,r}.
\]
The preceding bound on \(h_{z,r}\) implies
\[
|\mathbb E_r h_{z,r}|\le B,
\qquad
|h_{z,r}^{\circ}(S)|\le2B .
\]
For all sufficiently large \(r\), \(4M\le n_r\).  Applying \cref{obj:thm:exact-kneser-identity} to the centered table \(h_{z,r}^{\circ}\) and using the degree-\(\ge1\) spectral bound gives
\[
|C_{zz,r}|
\le
(M+1)\frac{2M}{n_r}(2B)^2 .
\]
Thus
\[
C_{11,r}\to0,
\qquad
C_{00,r}\to0 .
\]
Since \cref{obj:ass:stable-treatment-fraction} gives \(p_r\to p\in(0,1)\), also \(1-p_r\to1-p\in(0,1)\).  Hence
\[
\frac{C_{11,r}}{p_r}
+
\frac{C_{00,r}}{1-p_r}
\to0 .
\]
This proves the third displayed limit.
% lean: CausalSmith.Experimentation.DenseGroupPartitionProjectionPhase.armCrossCov_abs_le_ge_four

It remains to identify the variance limit.  The same stable-fraction assumption implies that, eventually,
\[
p_r>\frac p2,
\qquad
1-p_r>\frac{1-p}{2}.
\]
Together with \(G_r\to\infty\), this yields
\[
G_{1r}=p_rG_r\ge2,
\qquad
G_{0r}=(1-p_r)G_r\ge2
\]
for all sufficiently large \(r\).  On those rows, \cref{obj:thm:exact-pame-variance} gives
\[
\sigma_r^2
=
\frac{R_r}{G_r}
+
C_{\tau\tau,r}
-
\frac{C_{11,r}}{G_{1r}}
-
\frac{C_{00,r}}{G_{0r}} .
\]
Multiplying by \(G_r\), and using \(G_{1r}=p_rG_r\) and \(G_{0r}=(1-p_r)G_r\), gives the eventual identity
\[
G_r\sigma_r^2-(R_r-\rho E_{\tau,1,r})
=
\bigl(G_rC_{\tau\tau,r}+\rho E_{\tau,1,r}\bigr)
-
\left(
\frac{C_{11,r}}{p_r}
+
\frac{C_{00,r}}{1-p_r}
\right).
\]
Both terms on the right tend to zero by the preceding two steps.  Therefore
\[
G_r\sigma_r^2-(R_r-\rho E_{\tau,1,r})\to0,
\]
which is the first displayed limit.
% lean: CausalSmith.Experimentation.DenseGroupPartitionProjectionPhase.dense_projection_limit

Finally, the rowwise bound already proved gives
\[
\left|
G_r\sum_{k=2}^{M}\lambda_{r,k}\|\Pi_{r,k}d_r\|_r^2
\right|
\le
\frac{\Gamma}{n_r}
\le
\frac{C}{n_r},
\]
because \(C=\Gamma+1\) and \(n_r>0\).  Translating the row index \(r\) back to the theorem's notation gives the asserted bound for every row and completes the proof.
% lean: CausalSmith.Experimentation.DenseGroupPartitionProjectionPhase.dense_projection_limit
\end{proof}

\begin{proof}[Proof of \cref{obj:thm:cr2-phase-frontier}]
\begin{enumerate}
\item For \(z\in\{0,1\}\), set
\[
m_{1,n}:=G_{1n},\qquad m_{0,n}:=G_{0n}:=G_n-G_{1n},
\qquad
\pi_{1,n}:=p_n,\qquad \pi_{0,n}:=1-p_n .
\]
The dense schedule condition in \cref{obj:def:dense-schedule-class} gives \(G_n\to\infty\), \(p_n\to p\in(0,1)\), and hence
\[
m_{1,n}=p_nG_n\to\infty,
\qquad
m_{0,n}=(1-p_n)G_n\to\infty .
\]
Consequently \(m_{1,n}\ge2\) and \(m_{0,n}\ge2\) for all sufficiently large \(n\). Also, since \(MG_n\le n\) and \(M\ge2\), the population size tends to infinity, so \(4M\le n\) eventually. % lean: treated_tendsto_atTop, controls_tendsto_atTop

\item Let \(g_n:\Omega_{n,M}\to\mathbb R\) be any deterministic table with \(|g_n(S)|\le L\) for all \(S\in\Omega_{n,M}\), where \(L>0\). For \(\omega=(\mathbf A_n,\mathbf Z_n)\) in the support of the design in \cref{obj:def:random-partition-design}, define the arm-\(z\) selected mean by
\[
\overline g_{z,n}(\omega)
:=
\frac{1}{m_{z,n}}\sum_{j:Z_{jn}=z} g_n(A_{jn}),
\qquad
\mu_{g,n}:=\mathbb E_n g_n .
\]
The selected mean has expectation \(\mu_{g,n}\). Indeed, condition on any realized arm-position set of cardinality \(m_{z,n}\). Under \cref{obj:def:random-partition-design}, each selected group is uniform on \(\Omega_{n,M}\), and averaging over the independently chosen arm positions preserves this marginal law.

For the variance, define the centered table and the normalized disjointness operator by
\[
g_{n}^{\circ}(S):=g_n(S)-\mu_{g,n},
\qquad
(\mathsf K_n u)(S)
:=
\frac{1}{\#\{T\in\Omega_{n,M}:T\cap S=\varnothing\}}
\sum_{T\cap S=\varnothing}u(T),
\]
for \(S\in\Omega_{n,M}\). For all sufficiently large rows, \(m_{z,n}\ge2\) and \(2M\le n\). Conditional on the arm positions, two distinct selected groups form a uniform ordered disjoint pair, so
\[
\operatorname{Cov}\{g_n(A_{jn}),g_n(A_{\ell n})\}
=
\left\langle g_{n}^{\circ},\mathsf K_n g_{n}^{\circ}\right\rangle_n,
\qquad j\ne\ell .
\]
Expanding the square of the selected average gives
\[
\begin{aligned}
\operatorname{Var}\{\overline g_{z,n}\}
&=
\frac{1}{m_{z,n}^2}
\left\{
 m_{z,n}\operatorname{Var}_n(g_n)
 +m_{z,n}(m_{z,n}-1)
 \left\langle g_{n}^{\circ},\mathsf K_n g_{n}^{\circ}\right\rangle_n
\right\} \\
&=
\frac{\operatorname{Var}_n(g_n)}{m_{z,n}}
+
\left(1-\frac{1}{m_{z,n}}\right)
\left\langle g_{n}^{\circ},\mathsf K_n g_{n}^{\circ}\right\rangle_n .
\end{aligned}
\]

It remains to bound the covariance term. Since \(|g_n(S)|\le L\), uniform averaging gives \(|\mu_{g,n}|\le L\), hence \(|g_n^{\circ}(S)|\le2L\) and \(\|g_n^{\circ}\|_n^2\le(2L)^2\). Applying \cref{obj:thm:exact-kneser-identity} to the centered table gives
\[
\left\langle g_{n}^{\circ},\mathsf K_n g_{n}^{\circ}\right\rangle_n
=
\sum_{k=1}^{M}\lambda_{n,k}
\left\|\Pi_{n,k}g_{n}^{\circ}\right\|_n^2 .
\]
For rows with \(4M\le n\), the falling-factorial expression for \(\lambda_{n,k}\) satisfies, for \(1\le k\le M\),
\[
|\lambda_{n,k}|
\le
\left(\frac{2M}{n}\right)^k
\le
\frac{2M}{n} .
\]
The first inequality follows by bounding the numerator falling factorial by \(M^k\) and the denominator falling factorial by \((n/2)^k\); the second uses \(2M/n\le1\). Because each \(\Pi_{n,k}\) is an orthogonal projection,
\[
\left\|\Pi_{n,k}g_{n}^{\circ}\right\|_n^2
\le
\left\|g_{n}^{\circ}\right\|_n^2
\le
(2L)^2 .
\]
Therefore, for all sufficiently large \(n\),
\[
\left|
\left\langle g_{n}^{\circ},\mathsf K_n g_{n}^{\circ}\right\rangle_n
\right|
\le
(M+1)\frac{2M}{n}(2L)^2 .
\]
Since \(\operatorname{Var}_n(g_n)\le L^2\), \(m_{z,n}\to\infty\), and \(n\to\infty\), the displayed variance tends to zero. Chebyshev's inequality therefore gives
\[
\overline g_{z,n}-\mathbb E_n g_n\xrightarrow{p}0 .
\]
% lean: selectedMean_tendstoInProb

\item Apply the preceding conclusion first to \(g_n=h_{z,n}\), using \(|h_{z,n}|\le B\), and then to \(g_n=h_{z,n}^2\), using \(|h_{z,n}^2|\le B^2\). With
\[
\overline h_{z,n}(\omega):=\frac{1}{m_{z,n}}\sum_{j:Z_{jn}=z}h_{z,n}(A_{jn}),
\qquad
\overline {h^2}_{z,n}(\omega):=\frac{1}{m_{z,n}}\sum_{j:Z_{jn}=z}h_{z,n}(A_{jn})^2,
\]
we have
\[
\overline h_{z,n}-\mathbb E_n h_{z,n}\xrightarrow{p}0,
\qquad
\overline {h^2}_{z,n}-\mathbb E_n h_{z,n}^2\xrightarrow{p}0 .
\]
The sample-variance identity gives, eventually,
\[
s_{z,n}^2
=
\frac{m_{z,n}}{m_{z,n}-1}
\left\{
\overline {h^2}_{z,n}
-
\overline h_{z,n}^{\,2}
\right\}.
\]
Because \(m_{z,n}/(m_{z,n}-1)\to1\), and the selected means are bounded by \(B\), the square map preserves convergence in probability here. Thus
\[
s_{z,n}^2
\xrightarrow{p}
\mathbb E_n h_{z,n}^2-(\mathbb E_n h_{z,n})^2
=
V_{z,n}.
\]
% lean: armSampleVar_tendstoInProb

\item By \cref{obj:def:cr2-statistic} and \cref{obj:def:dense-correction},
\[
G_n\widehat V_{\mathrm{CR2},n}-R_n
=
\frac{s_{1,n}^2-V_{1,n}}{p_n}
+
\frac{s_{0,n}^2-V_{0,n}}{1-p_n}
\]
eventually. Since \(p_n\to p\in(0,1)\), both reciprocal factors are bounded deterministic sequences. Combining this boundedness with the two armwise convergences in the preceding step yields
\[
G_n\widehat V_{\mathrm{CR2},n}-R_n\xrightarrow{p}0 .
\]
% lean: cr2_centered_tendstoInProb

\item Define, for every \(n\),
\[
d_n:=R_n-\rho E_{\tau,1,n}.
\]
The first conclusion of \cref{obj:thm:dense-projection-limit} supplies
\[
G_n\sigma_n^2-d_n\to0 .
\]
Subtracting this deterministic convergence from the stochastic expansion just proved gives
\[
G_n\bigl(\widehat V_{\mathrm{CR2},n}-\sigma_n^2\bigr)
-\rho E_{\tau,1,n}
=
\bigl(G_n\widehat V_{\mathrm{CR2},n}-R_n\bigr)
-
\bigl(G_n\sigma_n^2-d_n\bigr)
\xrightarrow{p}0 .
\]
% lean: cr2_variance_gap_tendstoInProb

\item The scaled-variance nondegeneracy condition in \cref{obj:def:dense-schedule-class} gives
\[
\frac{c_\sigma}{2}<G_n\sigma_n^2
\]
eventually. Since \(G_n\sigma_n^2-d_n\to0\), for all sufficiently large \(n\),
\[
\left|G_n\sigma_n^2-d_n\right|<\frac{c_\sigma}{4},
\qquad\text{and hence}\qquad
\frac{c_\sigma}{4}<d_n .
\]
Thus \(d_n\) is eventually positive and bounded away from zero. % lean: scaledVariance_eventually_lower

\item Set
\[
r_n:=\frac{R_n}{d_n}.
\]
From the first expansion and the eventual lower bound on \(d_n\),
\[
\frac{G_n\widehat V_{\mathrm{CR2},n}}{d_n}-\frac{R_n}{d_n}
=
\bigl(G_n\widehat V_{\mathrm{CR2},n}-R_n\bigr)d_n^{-1}
\xrightarrow{p}0 .
\]
Similarly,
\[
\frac{G_n\sigma_n^2}{d_n}-1
=
\frac{G_n\sigma_n^2-d_n}{d_n}
\to0 .
\]
The deterministic target \(r_n\) is eventually bounded. To see this, bounded potential outcomes give \(V_{z,n}\le B^2\), while \(p_n\to p\in(0,1)\) gives
\[
\frac{p}{2}<p_n,
\qquad
\frac{1-p}{2}<1-p_n
\]
eventually. Hence, eventually,
\[
0\le R_n
\le
\frac{B^2}{p/2}+
\frac{B^2}{(1-p)/2}
=:\overline R,
\qquad
|r_n|\le \overline R\,\frac{4}{c_\sigma} .
\]
Let
\[
X_n(\omega):=\frac{G_n\widehat V_{\mathrm{CR2},n}(\omega)}{d_n},
\qquad
Y_n:=\frac{G_n\sigma_n^2}{d_n} .
\]
The preceding displays give \(X_n-r_n\xrightarrow{p}0\) and \(Y_n\to1\). Since \(Y_n\) is eventually bounded away from zero and \(r_n\) is eventually bounded,
\[
\frac{X_n}{Y_n}-r_n
=
\frac{(X_n-r_n)+r_n(1-Y_n)}{Y_n}
\xrightarrow{p}0 .
\]
Because \(X_n/Y_n=\widehat V_{\mathrm{CR2},n}/\sigma_n^2\) on the same eventual rows, this gives
\[
\frac{\widehat V_{\mathrm{CR2},n}}{\sigma_n^2}
\xrightarrow{p}
r_n
=
\frac{R_n}{R_n-\rho E_{\tau,1,n}} .
\]
% lean: cr2_ratio_leading_tendstoInProb

\item Define
\[
q_n:=\frac{\rho E_{\tau,1,n}}{d_n}
=
\frac{\rho E_{\tau,1,n}}{R_n-\rho E_{\tau,1,n}} .
\]
For all sufficiently large \(n\), \(d_n\ne0\), and hence
\[
r_n
=
\frac{R_n}{d_n}
=
\frac{d_n+\rho E_{\tau,1,n}}{d_n}
=
1+q_n .
\]
Moreover \(\rho\ge0\) by the sampling-fraction condition, and \(E_{\tau,1,n}\ge0\) because it is a squared norm. Since \(d_n>0\) eventually,
\[
r_n=1+q_n\ge1
\]
eventually. % lean: haq, haOne

\item Fix \(\varepsilon>0\). On every row for which \(r_n\ge1\),
\[
\left\{
\frac{\widehat V_{\mathrm{CR2},n}}{\sigma_n^2}<1-\varepsilon
\right\}
\subseteq
\left\{
\varepsilon\le
\left|
\frac{\widehat V_{\mathrm{CR2},n}}{\sigma_n^2}-r_n
\right|
\right\}.
\]
The right-hand probability tends to zero because
\(\widehat V_{\mathrm{CR2},n}/\sigma_n^2\xrightarrow{p}r_n\). Therefore
\[
\Pr\!\left\{
\widehat V_{\mathrm{CR2},n}/\sigma_n^2<1-\varepsilon
\right\}\to0 .
\]
% lean: lower_tail_vanishes_of_tendstoInProb

\item It remains to prove the equivalence. Suppose first that
\[
\frac{\widehat V_{\mathrm{CR2},n}}{\sigma_n^2}\xrightarrow{p}1 .
\]
Together with
\[
\frac{\widehat V_{\mathrm{CR2},n}}{\sigma_n^2}\xrightarrow{p}r_n,
\]
uniqueness of deterministic probability limits gives \(r_n\to1\): indeed, if \(|r_n-1|\ge\eta\), the triangle inequality forces at least one of
\[
\left|
\frac{\widehat V_{\mathrm{CR2},n}}{\sigma_n^2}-r_n
\right|\ge\frac{\eta}{2},
\qquad
\left|
\frac{\widehat V_{\mathrm{CR2},n}}{\sigma_n^2}-1
\right|\ge\frac{\eta}{2}
\]
to occur. Since \(q_n=r_n-1\) eventually, this yields \(q_n\to0\), namely
\[
\frac{\rho E_{\tau,1,n}}{R_n-\rho E_{\tau,1,n}}\to0 .
\]
% lean: tendstoInProb_target_unique

\item Conversely, assume
\[
\frac{\rho E_{\tau,1,n}}{R_n-\rho E_{\tau,1,n}}\to0 .
\]
That is, \(q_n\to0\). Since \(r_n=1+q_n\) eventually, \(r_n\to1\). Retargeting the already proved convergence
\[
\frac{\widehat V_{\mathrm{CR2},n}}{\sigma_n^2}\xrightarrow{p}r_n
\]
to the deterministic limit \(1\) gives
\[
\frac{\widehat V_{\mathrm{CR2},n}}{\sigma_n^2}\xrightarrow{p}1 .
\]
This proves the asserted if and only if. % lean: FiniteDesign.TendstoInProb.retarget
\end{enumerate}
\end{proof}

\begin{proof}[Proof of \cref{obj:thm:sparse-beyond-birthday}]
Since \(M\ge 2\), in particular \(M>0\). We prove the three asserted conclusions in the order in which they are stated.
\begin{enumerate}
\item \textbf{Sparse dense arrays.}
Let \(A\) be a feasible triangular schedule array with common group size \(M\), and suppose that \(A\in\mathcal C_{\mathrm{dense}}\) in \cref{obj:def:dense-schedule-class} with limiting grouped-unit fraction \(\rho=0\). Equip each row with the canonical Johnson projections. The canonical decomposition and Kneser-spectrum hypotheses required by \cref{obj:thm:cr2-phase-frontier} are supplied by \cref{obj:lem:classical-johnson-decomposition,obj:lem:classical-kneser-spectrum}. Applying \cref{obj:thm:cr2-phase-frontier} with \(\rho=0\), its ratio criterion becomes
\[
\frac{0\cdot E_{\tau,1,n}}{R_n-0\cdot E_{\tau,1,n}}=0
\qquad\text{for every row }n,
\]
so the criterion converges to \(0\). Hence
\[
\frac{\widehat V_{\mathrm{CR2},n}}{\sigma_n^2}
\xrightarrow{p}1 .
\]
% lean: sparse_dense_ratio_consistency

\item \textbf{Birthday counts.}
Define
\[
G_n^{\mathrm{bd}}
=
\left\lfloor \frac{n^{3/4}}{M}\right\rfloor,
\qquad
N_n^{\mathrm{bd}}=M G_n^{\mathrm{bd}} .
\]
For \(n\ge 1\),
\[
N_n^{\mathrm{bd}}
=
M\left\lfloor \frac{n^{3/4}}{M}\right\rfloor
\le
M\frac{n^{3/4}}{M}
=
n^{3/4}
\le n,
\]
so \(N_n^{\mathrm{bd}}\le n\) eventually. For \(n>0\),
\[
\frac{N_n^{\mathrm{bd}}}{n}
=
\frac{\left\lfloor n^{3/4}/M\right\rfloor}{n^{3/4}/M}\,
n^{-1/4}.
\]
Since \(n^{3/4}/M\to\infty\), the floor ratio tends to \(1\), while \(n^{-1/4}\to0\). Therefore
\[
\frac{N_n^{\mathrm{bd}}}{n}\to0.
\]
For divergence of the squared count, the elementary floor bound gives
\[
n^{3/4}-M
\le
M\left\lfloor \frac{n^{3/4}}{M}\right\rfloor
=
N_n^{\mathrm{bd}} .
\]
Eventually \(n^{3/4}\ge 2M\), and hence
\[
\frac{n^{3/4}}{2}\le N_n^{\mathrm{bd}}.
\]
Consequently, for all sufficiently large \(n\),
\[
\frac{(N_n^{\mathrm{bd}})^2}{n}
\ge
\frac{(n^{3/4}/2)^2}{n}
=
\frac14 n^{1/2}
\to\infty .
\]
% lean: birthdayBenchmark_proof

\item \textbf{Birthday-aligned arrays.}
Let \(A\) be a feasible triangular schedule array satisfying the rowwise birthday alignment, \cref{obj:ass:group-count-growth,obj:ass:stable-treatment-fraction,obj:ass:bounded-potential-outcomes,obj:ass:scaled-variance-nondegeneracy}. For row \(r\in\mathbb N\), write \(G_r\) for the group count, \(n_r\) for the population size, and \(N_r=MG_r\) for the grouped-unit count. From \cref{obj:ass:group-count-growth}, \(G_r\to\infty\), and feasibility gives
\[
G_r\le M G_r=N_r\le n_r .
\]
Thus \(n_r\to\infty\). The alignment condition says
\[
N_r
=
M\left\lfloor \frac{n_r^{3/4}}{M}\right\rfloor
=
N_{n_r}^{\mathrm{bd}},
\]
and the birthday-count limit from the previous step therefore yields
\[
\frac{N_r}{n_r}
=
\frac{N_{n_r}^{\mathrm{bd}}}{n_r}
\to 0 .
\]
Together with \(0\le 0\le 1\), this is exactly \cref{obj:ass:sampling-fraction} with limiting grouped-unit fraction \(\rho=0\). The assumed group-count growth, stable treatment fraction, bounded potential outcomes, and scaled-variance nondegeneracy then place \(A\) in \(\mathcal C_{\mathrm{dense}}\) with \(\rho=0\) by \cref{obj:def:dense-schedule-class}. The first step applies and gives
\[
\frac{\widehat V_{\mathrm{CR2},n}}{\sigma_n^2}
\xrightarrow{p}1 .
\]
% lean: sparse_beyond_birthday
\end{enumerate}
\end{proof}

\begin{proof}[Proof of \cref{obj:prop:clubsandwich-consumer}]
\begin{enumerate}
\item
Write \(T_n\subseteq\{1,\ldots,G_n\}\) for the realized treated-group set, so
\(|T_n|=G_{1n}\), and write \(C_n=T_n^{c}\), so
\(|C_n|=G_{0n}=G_n-G_{1n}\).  For a group \(g\), define
\[
G_{g,n}:=
\begin{cases}
G_{1n},& g\in T_n,\\
G_{0n},& g\in C_n.
\end{cases}
\]
The assumptions \(G_{1n}\ge2\) and \(G_{0n}\ge2\) imply \(G_{1n}\le G_n\),
\(G_n>0\), \(M>0\), \(G_{0n}>0\), and the real denominators
\(G_{1n}\) and \(G_{0n}\) are nonzero.  The two versioned source identities give
\[
{\tt sandwich::bread}({\tt fit})
=
G_nM\,(X^{\mathsf T}X)^{-1}
\]
for the unweighted full-rank least-squares fit
\citep[R/bread.R, lines 10--15]{Sandwich2026}, and
\[
\Sigma
=
\widetilde B
\left\{
\sum_{g=1}^{G_n}X_g^{\mathsf T}A_ge_ge_g^{\mathsf T}A_g^{\mathsf T}X_g
\right\}
\widetilde B
\]
for the CR2 call with omitted \({\tt target}\), \({\tt inverse\_var}\), and
\({\tt form}\)
\citep[R/lm.R, lines 47--52 and 70--72; R/S3-methods.R, lines 21--33 and
63--65; R/clubSandwich.R, lines 167--175, 216--225, and 237--287;
R/CR-adjustments.R, lines 5--7 and 22--42]{ClubSandwich2026}.  The stated bread
scaling therefore yields
\[
\widetilde B=(X^{\mathsf T}X)^{-1}.
\]
% lean: clubsandwich_cr2_consumer_identity

\item
The fitted design has group block
\[
X_g=\mathbf 1_M x_g^{\mathsf T},
\qquad
x_g=
\begin{cases}
(1,1)^{\mathsf T},& g\in T_n,\\
(1,0)^{\mathsf T},& g\in C_n,
\end{cases}
\]
where \(\mathbf 1_M\) is the \(M\)-vector of ones.  Summing the block
cross-products gives
\[
X^{\mathsf T}X
=
\sum_{g=1}^{G_n}X_g^{\mathsf T}X_g
=
M
\begin{pmatrix}
G_n&G_{1n}\\
G_{1n}&G_{1n}
\end{pmatrix}.
\]
Since the determinant is \(M^2G_{1n}G_{0n}>0\), inversion gives the entries used
below:
\[
\widetilde B_{10}=\widetilde B_{01}
=-\frac{1}{MG_{0n}},
\qquad
\widetilde B_{11}
=
\frac{1}{MG_{0n}}+\frac{1}{MG_{1n}},
\qquad
\widetilde B_{00}
=
\frac{1}{MG_{0n}}.
\]
% lean: clubsandwich_cr2_consumer_identity

\item
For every group \(g\), the leverage block is
\[
H_g=X_g\widetilde B X_g^{\mathsf T}.
\]
Using the preceding inverse entries and the two possible values of \(x_g\),
\[
(H_g)_{ij}
=
\frac{1}{M G_{g,n}},
\qquad
1\le i,j\le M.
\]
Thus
\[
I_M-H_g
=
I_M-\frac{1}{G_{g,n}}P_M,
\qquad
P_M:=\frac{1}{M}\mathbf 1_M\mathbf 1_M^{\mathsf T}.
\]
Let \(q=G_{g,n}\).  Since \(q\ge2\), the symmetric positive-definite matrix
\(A_g\) satisfying \(A_gA_g(I_M-H_g)=I_M\) is determined by its action on the
span of \(\mathbf 1_M\) and on its orthogonal complement.  For
\[
u:=\mathbf 1_M,
\qquad
v_j:=e_j-\frac1M u,
\]
one has
\[
(I_M-H_g)v_j=v_j,
\qquad
(I_M-H_g)u=\frac{q-1}{q}u.
\]
The identity \(A_gA_g(I_M-H_g)=I_M\), together with positive definiteness,
therefore gives
\[
A_gv_j=v_j,
\qquad
A_gu=\sqrt{\frac{q}{q-1}}\,u.
\]
Combining these two relations column by column,
\[
(A_g)_{ij}
=
\mathbf 1\{i=j\}
+
\frac{\sqrt{G_{g,n}/(G_{g,n}-1)}-1}{M}.
\]
% lean: posDef_rankOne_inverse_sqrt

\item
Summing the displayed entries down any column gives
\[
\sum_{i=1}^{M}(A_g)_{ij}
=
\sqrt{\frac{G_{g,n}}{G_{g,n}-1}}.
\]
Consequently
\[
\mathbf 1_M^{\mathsf T}A_ge_g
=
\sum_{i=1}^{M}(A_ge_g)_i
=
\sqrt{\frac{G_{g,n}}{G_{g,n}-1}}
\sum_{j=1}^{M}e_{gj}.
\]
The residual-mean condition in the statement says
\[
\frac1M\sum_{j=1}^{M}e_{gj}
=
H_{g,n}^{\mathrm{obs}}-\bar H_{z_g,n},
\]
where
\[
H_{g,n}^{\mathrm{obs}}:=\operatorname{obsGroupMean}(Y,\mathbf A_n,\mathbf Z_n;g),
\qquad
\bar H_{z,n}:=\operatorname{armObsMean}(Y,\mathbf A_n,\mathbf Z_n;z),
\]
and \(z_g=1\) on \(T_n\), \(z_g=0\) on \(C_n\).  Hence
\[
\sum_{j=1}^{M}e_{gj}
=
M\{H_{g,n}^{\mathrm{obs}}-\bar H_{z_g,n}\}.
\]
% lean: clubsandwich_cr2_consumer_identity

\item
Define the treatment-coordinate adjusted cluster score by
\[
q_{g,n}
:=
\left[
\widetilde B\,X_g^{\mathsf T}A_ge_g
\right]_1 .
\]
Using \(X_g=\mathbf 1_Mx_g^{\mathsf T}\), the inverse entries in Step 2, and the
residual sum from Step 4 gives the two cases
\[
q_{g,n}
=
\begin{cases}
\displaystyle
\sqrt{\frac{G_{1n}}{G_{1n}-1}}\,
\frac{H_{g,n}^{\mathrm{obs}}-\bar H_{1,n}}{G_{1n}},
& g\in T_n,\\[1.2em]
\displaystyle
-\sqrt{\frac{G_{0n}}{G_{0n}-1}}\,
\frac{H_{g,n}^{\mathrm{obs}}-\bar H_{0,n}}{G_{0n}},
& g\in C_n.
\end{cases}
\]
% lean: clubsandwich_cr2_consumer_identity

\item
Because \(\widetilde B\) is symmetric, each summand in the CR2 sandwich is a
rank-one outer product after premultiplication and postmultiplication by
\(\widetilde B\).  Indeed, with
\[
P_g:=\widetilde B X_g^{\mathsf T}A_g,
\]
\[
\widetilde B X_g^{\mathsf T}A_ge_ge_g^{\mathsf T}A_g^{\mathsf T}X_g\widetilde B
=
P_ge_ge_g^{\mathsf T}P_g^{\mathsf T}
=
(P_ge_g)(P_ge_g)^{\mathsf T}.
\]
Therefore
\[
\Sigma
=
\sum_{g=1}^{G_n}
\left(\widetilde B X_g^{\mathsf T}A_ge_g\right)
\left(\widetilde B X_g^{\mathsf T}A_ge_g\right)^{\mathsf T},
\]
and its treatment-coordinate entry is
\[
\Sigma_{11}
=
\sum_{g=1}^{G_n}q_{g,n}^{\,2}.
\]
% lean: matrix_sandwich_rank_one

\item
Squaring the two cases in Step 5 and using \(G_{1n},G_{0n}\ge2\) gives
\[
\Sigma_{11}
=
\sum_{g\in T_n}
\frac{(H_{g,n}^{\mathrm{obs}}-\bar H_{1,n})^2}{(G_{1n}-1)G_{1n}}
+
\sum_{g\in C_n}
\frac{(H_{g,n}^{\mathrm{obs}}-\bar H_{0,n})^2}{(G_{0n}-1)G_{0n}}.
\]
Equivalently,
\[
\Sigma_{11}
=
\frac{1}{G_{1n}}
\frac{\sum_{g\in T_n}(H_{g,n}^{\mathrm{obs}}-\bar H_{1,n})^2}{G_{1n}-1}
+
\frac{1}{G_{0n}}
\frac{\sum_{g\in C_n}(H_{g,n}^{\mathrm{obs}}-\bar H_{0,n})^2}{G_{0n}-1}.
\]
By \cref{obj:def:cr2-statistic}, the two displayed fractions are
\(s_{1,n}^2\) and \(s_{0,n}^2\), respectively, and hence
\[
\Sigma_{11}
=
\frac{s_{1,n}^2}{G_{1n}}
+
\frac{s_{0,n}^2}{G_{0n}}
=
\widehat V_{\mathrm{CR2},n}.
\]
% lean: clubsandwich_cr2_consumer_identity
\end{enumerate}
\end{proof}

\begin{proof}[Proof of \cref{obj:prop:eight-unit-witness}]
Let
\[
\Sigma_8:\Omega_{8,2}\to\mathbb R,
\qquad
\Sigma_8(A)=\sum_{u\in A} a_u .
\]
By \cref{obj:def:eight-unit-witness,obj:def:group-table},
\[
h_{1,8}(A)=\frac{\Sigma_8(A)}{2},
\qquad
h_{0,8}(A)=0
\quad (A\in\Omega_{8,2}).
\]
The vector \(\mathbf a_8\) has four \(+1\) entries and four \(-1\) entries, so its population mean is zero and its population variance is one. The usual fixed-size sample-mean identity for a uniform two-subset gives
\[
\mathbb E_8 h_{1,8}=0,
\qquad
V_{1,8}
=
\operatorname{Var}_{A\sim\mathrm{Unif}(\Omega_{8,2})}
  \left\{\frac{\Sigma_8(A)}{2}\right\}
=
\frac{8-2}{2(8-1)}\cdot 1
=
\frac{3}{7}.
\]
Since \(h_{0,8}\) is identically zero, its centered table is zero, hence
\[
V_{0,8}=0,
\qquad
C_{00,8}=C_{10,8}=C_{01,8}=0 .
\]
% lean: witness8_armVar_treated

For the treated ordered-disjoint covariance, write the eight-unit population and its two-subset slice as
\[
\mathcal U_8=\{1,\ldots,8\},
\qquad
\Omega_{8,2}=\{A\subseteq\mathcal U_8: |A|=2\}.
\]
Use ordered pairs
\[
(A,C)\in\Omega_{8,2}\times\Omega_{8,2},
\qquad
A\cap C=\varnothing .
\]
There are
\[
\binom82\binom62=420
\]
such ordered pairs. For fixed \(A\), each unit in the six-point complement of \(A\) appears in five two-subsets of that complement, and \(\sum_{u\in\mathcal U_8}a_u=0\). Therefore
\[
\sum_{\substack{C\in\Omega_{8,2}\\ C\cap A=\varnothing}}\Sigma_8(C)
=
5\sum_{u\notin A}a_u
=
-5\Sigma_8(A).
\]
Also \(\sum_{A\in\Omega_{8,2}}\Sigma_8(A)^2=48\), because the six same-positive pairs and six same-negative pairs contribute \(4\) each and the sixteen mixed pairs contribute \(0\). Thus
\[
\sum_{\substack{A,C\in\Omega_{8,2}\\ A\cap C=\varnothing}}
\Sigma_8(A)\Sigma_8(C)
=
-5\sum_{A\in\Omega_{8,2}}\Sigma_8(A)^2
=
-240.
\]
Using \(\mathbb E_8h_{1,8}=0\), this gives
\[
C_{11,8}
=
\frac1{420}
\sum_{\substack{A,C\in\Omega_{8,2}\\ A\cap C=\varnothing}}
\frac{\Sigma_8(A)}2\frac{\Sigma_8(C)}2
=
\frac{-240}{1680}
=
-\frac17 .
\]
Consequently
\[
C_{\tau\tau,8}
=
C_{11,8}+C_{00,8}-2C_{10,8}
=
-\frac17 .
\]
% lean: witness8_crossCov_treated

Define the arm-difference table
\[
\Delta_8:\Omega_{8,2}\to\mathbb R,
\qquad
\Delta_8(A)=h_{1,8}(A)-h_{0,8}(A).
\]
Since \(h_{0,8}=0\), \(\Delta_8=h_{1,8}\), so
\[
\mathbb E_8\Delta_8=0,
\qquad
\operatorname{Var}_{\mathbb E_8}(\Delta_8)=\frac37 .
\]
Set
\[
\eta_1=\|\Pi_{8,1}\Delta_8\|_8^2,
\qquad
\eta_2=\|\Pi_{8,2}\Delta_8\|_8^2 .
\]
The centered Johnson decomposition assumed in the proposition gives
\[
\Delta_8=\Pi_{8,1}\Delta_8+\Pi_{8,2}\Delta_8
\]
and the two summands are orthogonal. Hence
\[
\eta_1+\eta_2
=
\|\Delta_8\|_8^2
=
\frac37 .
\]
By \cref{obj:thm:exact-kneser-identity}, instantiated with \(n=8\) and \(M=2\),
\[
C_{\tau\tau,8}
=
\lambda_{8,1}\eta_1+\lambda_{8,2}\eta_2 .
\]
The Kneser eigenvalues are
\[
\lambda_{8,1}=-\frac{2}{6}=-\frac13,
\qquad
\lambda_{8,2}=\frac{2\cdot 1}{6\cdot 5}=\frac1{15}.
\]
Using \(C_{\tau\tau,8}=-1/7\),
\[
-\frac13\eta_1+\frac1{15}\eta_2=-\frac17 .
\]
Together with \(\eta_1+\eta_2=3/7\), this linear system gives
\[
\eta_1=\frac37 .
\]
By \cref{obj:def:dense-correction}, \(\eta_1=E_{\tau,1,8}\), so
\[
E_{\tau,1,8}=\frac37 .
\]
% lean: witness8_degreeOne_from_spectrum

It remains to evaluate the design variance and the CR2 expectation. Here
\[
G_8=4,\qquad G_{1,8}=2,\qquad G_{0,8}=2,\qquad p_8=\frac12 .
\]
By \cref{obj:def:dense-correction},
\[
R_8=\frac{V_{1,8}}{p_8}+\frac{V_{0,8}}{1-p_8}
=
\frac{3/7}{1/2}+0
=
\frac67 .
\]
Applying \cref{obj:thm:exact-pame-variance} to the feasible eight-unit design gives
\[
\sigma_8^2
=
\frac{R_8}{G_8}
+
C_{\tau\tau,8}
-
\frac{C_{11,8}}{G_{1,8}}
-
\frac{C_{00,8}}{G_{0,8}}.
\]
Substituting the displayed values,
\[
\sigma_8^2
=
\frac{6/7}{4}
-\frac17
-\frac{-1/7}{2}
-0
=
\frac17 .
\]
The same result in \cref{obj:thm:exact-pame-variance}, now for the statistic in \cref{obj:def:cr2-statistic}, gives
\[
\mathbb E\!\left[\widehat V_{\mathrm{CR2},8}\right]
=
\frac{V_{1,8}-C_{11,8}}{G_{1,8}}
+
\frac{V_{0,8}-C_{00,8}}{G_{0,8}}.
\]
Therefore
\[
\mathbb E\!\left[\widehat V_{\mathrm{CR2},8}\right]
=
\frac{3/7+1/7}{2}+0
=
\frac27 .
\]
% lean: eight_unit_witness_moments
\end{proof}

\begin{proof}[Proof of \cref{obj:prop:rademacher-mixture-separation}]
\begin{enumerate}
\item Fix a row \(n\) and a realized design \(w=(\mathbf A_n,\mathbf Z_n)\).  Under the independent-sign prior, define the selected sign field
\[
\xi_i^{w}:=\xi_{i,\alpha_w(i)},
\qquad
\alpha_w(i):=\mathbf 1\{i\text{ is observed in a treated realized group under }w\}.
\]
Because the realized groups are pairwise disjoint, every observed unit has \(\alpha_w(i)\) equal to its realized treatment arm.  Hence the observation produced from the independent schedule is exactly the observation produced from the common-sign schedule with signs \((\xi_i^w)_{i\in\mathcal U_n}\).  For fixed \(w\), the vector \((\xi_i^w)_i\) has the same product Rademacher law as the common-sign vector.  Finite Fubini then gives, for every statistic \(\varphi\),
\[
\mathbb E_{\mathbb H_n^{\mathrm{same}}}\mathbb E_{\mathbf Z_n}\{\varphi(\mathcal O_n)\}
=
\mathbb E_{\mathbf Z_n}\mathbb E_{\mathbb H_n^{\mathrm{same}}}\{\varphi(\mathcal O_n)\}
=
\mathbb E_{\mathbf Z_n}\mathbb E_{\mathbb H_n^{\mathrm{ind}}}\{\varphi(\mathcal O_n)\}
=
\mathbb E_{\mathbb H_n^{\mathrm{ind}}}\mathbb E_{\mathbf Z_n}\{\varphi(\mathcal O_n)\}.
\]
% lean: CausalSmith.Experimentation.DenseGroupPartitionProjectionPhase.sameObservedExpectation_eq_independent

\item For any deterministic vector \(\boldsymbol x=(x_i)_{i\in\mathcal U_n}\), write
\[
\bar x_n:=\frac1n\sum_{i\in\mathcal U_n}x_i,
\qquad
\operatorname{PV}_n(\boldsymbol x):=\frac1{n-1}\sum_{i\in\mathcal U_n}(x_i-\bar x_n)^2,
\qquad
m_{\boldsymbol x}(Q):=\frac1M\sum_{i\in Q}x_i
\]
for \(Q\in\Omega_{n,M}\).  Sampling without replacement gives
\[
\operatorname{Var}_{Q\sim\mathrm{Unif}(\Omega_{n,M})}\{m_{\boldsymbol x}(Q)\}
=
\left(\frac1M-\frac1n\right)\operatorname{PV}_n(\boldsymbol x).
\]
If \(x_i\in\{-1,1\}\), then
\[
\operatorname{PV}_n(\boldsymbol x)
=
\frac n{n-1}\left(1-\bar x_n^2\right).
\]
Under the common-sign prior, \(\mathbb E\bar\xi_n=0\) and
\[
\operatorname{Var}(\bar\xi_n)=\frac1n,
\]
so \(\bar\xi_n\to 0\) in probability.  Since \cref{obj:ass:group-count-growth} implies \(n\to\infty\),
\[
\left(\frac1M-\frac1n\right)\frac n{n-1}\to\frac1M,
\]
and therefore, for each arm \(z\),
\[
V_{z,n}\stackrel{p}{\longrightarrow}\frac1M
\qquad
\text{under }\mathbb H_n^{\mathrm{same}}.
\]
% lean: CausalSmith.Experimentation.DenseGroupPartitionProjectionPhase.scheduleArray_same_armVar_tendstoInProb

\item The same calculation applies arm by arm under \(\mathbb H_n^{\mathrm{ind}}\).  For each fixed arm \(z\),
\[
\bar\xi_{z,n}:=\frac1n\sum_{i\in\mathcal U_n}\xi_{i,z},
\qquad
\mathbb E\bar\xi_{z,n}=0,
\qquad
\operatorname{Var}(\bar\xi_{z,n})=\frac1n,
\]
so \(\bar\xi_{z,n}\to0\) in probability and
\[
V_{z,n}\stackrel{p}{\longrightarrow}\frac1M
\qquad
\text{under }\mathbb H_n^{\mathrm{ind}}.
\]
% lean: CausalSmith.Experimentation.DenseGroupPartitionProjectionPhase.scheduleArray_independent_armVar_tendstoInProb

\item Under \(\mathbb H_n^{\mathrm{same}}\), the two arm tables coincide pointwise:
\[
h_{1,n}(Q)-h_{0,n}(Q)=0
\qquad
(Q\in\Omega_{n,M}).
\]
By linearity of the Johnson projection,
\[
E_{\tau,1,n}
=
\left\|\Pi_{n,1}(h_{1,n}-h_{0,n})\right\|_n^2
=
\|\Pi_{n,1}0\|_n^2
=
0.
\]
Thus \(E_{\tau,1,n}\stackrel{p}{\to}0\) under \(\mathbb H_n^{\mathrm{same}}\).
% lean: CausalSmith.Experimentation.DenseGroupPartitionProjectionPhase.scheduleArray_same_degreeOne_tendstoInProb

\item Under \(\mathbb H_n^{\mathrm{ind}}\), put
\[
d_i:=\xi_{i,1}-\xi_{i,0},
\qquad
\bar d_n:=\frac1n\sum_i d_i,
\qquad
\bar c_n:=\frac1n\sum_i \xi_{i,1}\xi_{i,0}.
\]
The arm-difference table is the slice sample mean
\[
h_{1,n}(Q)-h_{0,n}(Q)=m_{\boldsymbol d}(Q).
\]
This function is a linear combination of degree-one inclusion indicators; after subtracting its slice mean, it is orthogonal to constants and lies in \(\mathcal H_{n,1}\).  Hence the degree-one projection equals the centered sample mean, and
\[
E_{\tau,1,n}
=
\left(\frac1M-\frac1n\right)\operatorname{PV}_n(\boldsymbol d).
\]
Since
\[
\operatorname{PV}_n(\boldsymbol d)
=
\frac n{n-1}
\left(2-2\bar c_n-\bar d_n^2\right),
\]
while \(\bar c_n\to0\), \(\bar\xi_{1,n}\to0\), and \(\bar\xi_{0,n}\to0\) in probability, we get
\[
E_{\tau,1,n}\stackrel{p}{\longrightarrow}\frac2M
\qquad
\text{under }\mathbb H_n^{\mathrm{ind}}.
\]
% lean: CausalSmith.Experimentation.DenseGroupPartitionProjectionPhase.scheduleArray_independent_degreeOne_tendstoInProb

\item By \cref{obj:ass:stable-treatment-fraction},
\[
p_n\to p,\qquad 0<p<1,
\]
and hence
\[
\frac1{p_n}\to\frac1p,
\qquad
\frac1{1-p_n}\to\frac1{1-p}.
\]
Consequently, for either prior,
\[
\Theta_n:=
\frac{V_{1,n}}{p_n}+\frac{V_{0,n}}{1-p_n}
\stackrel{p}{\longrightarrow}
\frac1M\left(\frac1p+\frac1{1-p}\right)
=
\frac1{M p(1-p)}.
\]
% lean: CausalSmith.Experimentation.DenseGroupPartitionProjectionPhase.rademacher_mixture_separation

\item Since \(0<G_{1n}<G_n\), every row has \(G_n\ge2\), and therefore \(2M\le M G_n\le n\).  The degree-one Kneser eigenvalue in \cref{obj:thm:exact-kneser-identity} is
\[
\lambda_{n,1}=-\frac{M}{n-M}.
\]
Thus \(\lambda_{n,1}\to0\).  Writing \(N_n=M G_n\) for the grouped-unit count and using \cref{obj:ass:sampling-fraction},
\[
G_n\lambda_{n,1}
=
-\frac{M G_n}{n-M}
=
-\frac{N_n/n}{1-M/n}
\longrightarrow
-\rho .
\]
% lean: CausalSmith.Experimentation.DenseGroupPartitionProjectionPhase.kneserEigenvalue_one

\item The exact rowwise variance identity is used on the eventual rows where both treatment arms contain at least two realized groups.  These rows are eventual because \(p_n\to p\in(0,1)\) and \(G_n\to\infty\): eventually \(p_n>p/2\) and \(G_n>4/p\), forcing \(G_{1n}\ge2\), and eventually \(p_n<(1+p)/2\) and \(G_n>4/(1-p)\), forcing \(G_{0n}\ge2\).

On those rows, \cref{obj:thm:exact-pame-variance,obj:thm:exact-kneser-identity} give
\[
G_n\sigma_n^2
=
\Theta_n+G_n\lambda_{n,1}E_{\tau,1,n}
-\lambda_{n,1}\Theta_n
=
(1-\lambda_{n,1})\Theta_n
+
G_n\lambda_{n,1}E_{\tau,1,n}.
\]
Combining the convergence statements above yields
\[
G_n\sigma_n^2
\stackrel{p}{\longrightarrow}
\frac1{M p(1-p)}
\qquad
\text{under }\mathbb H_n^{\mathrm{same}},
\]
and
\[
G_n\sigma_n^2
\stackrel{p}{\longrightarrow}
\frac1{M p(1-p)}-\frac{2\rho}{M}
\qquad
\text{under }\mathbb H_n^{\mathrm{ind}}.
\]
% lean: CausalSmith.Experimentation.DenseGroupPartitionProjectionPhase.scaledSigmaSq_eq_of_additive_armTables

\item The separation is positive because \(M>0\) and \(\rho>0\):
\[
0<\frac{2\rho}{M}.
\]
Also \cref{obj:ass:sampling-fraction} gives \(\rho\le1\), and \(p(1-p)\le1/4\) follows from \((p-1/2)^2\ge0\).  Since \(p(1-p)>0\),
\[
\frac{1}{M p(1-p)}-\frac{2\rho}{M}
=
\frac1M\left(\frac1{p(1-p)}-2\rho\right)
\ge
\frac1M(4-2\rho)
\ge
\frac2M .
\]
% lean: CausalSmith.Experimentation.DenseGroupPartitionProjectionPhase.rademacher_mixture_separation

\item Finally, if \(X_n\to c\) in probability and \(t<c\), then
\[
\Pr\{X_n<t\}
\le
\Pr\{|X_n-c|>c-t\}
\to0,
\qquad
\Pr\{X_n\ge t\}\to1.
\]
Apply this with \(X_n=G_n\sigma_n^2\).  The same-sign limit exceeds the independent-sign limit by \(2\rho/M\), so every
\[
0<c_\sigma<\frac{1}{M p(1-p)}-\frac{2\rho}{M}
\]
lies below both limits.  Therefore
\[
\Pr_{\mathbb H_n^{\mathrm{same}}}\{G_n\sigma_n^2\ge c_\sigma\}\to1,
\qquad
\Pr_{\mathbb H_n^{\mathrm{ind}}}\{G_n\sigma_n^2\ge c_\sigma\}\to1.
\]
% lean: CausalSmith.Experimentation.DenseGroupPartitionProjectionPhase.probability_ge_tendsto_one_of_tendstoInProb
\end{enumerate}
\end{proof}

\begin{proof}[Proof of \cref{obj:thm:qv-diagonal-impossibility}]
\begin{enumerate}
\item Define
\[
d_{\mathrm{same}}=\frac{1}{M p(1-p)},
\qquad
d_{\mathrm{ind}}=d_{\mathrm{same}}-\frac{2\rho}{M}.
\]
By \cref{obj:prop:rademacher-mixture-separation}, the unconditioned same-sign and independent-arm one-realization mixtures have identical bounded-test expectations, and the scaled exact variances satisfy
\[
G_n\sigma_n^2(Y_n^{\mathrm{same}})\xrightarrow{p}d_{\mathrm{same}},
\qquad
G_n\sigma_n^2(Y_n^{\mathrm{ind}})\xrightarrow{p}d_{\mathrm{ind}}
\]
under their respective Rademacher schedule priors. The same result also gives \(0<2\rho/M\), hence
\[
d_{\mathrm{ind}}<d_{\mathrm{same}}.
\]
The hypothesis on \(c_\sigma\) gives \(c_\sigma<d_{\mathrm{ind}}\), and therefore also \(c_\sigma<d_{\mathrm{same}}\). % lean: rademacher_mixture_separation

\item We extract deterministic high-probability supports from these two convergence-in-probability statements. For \(j\in\{\mathrm{same},\mathrm{ind}\}\), let \(d_j\) denote the corresponding limit and let
\[
X_{j,n}(\omega)=G_n\sigma_n^2(Y_n^{j,\omega})
\]
for a schedule-prior sign configuration \(\omega\). Since \(X_{j,n}\to d_j\) in probability, for each integer \(k\ge1\) choose \(N_{j,k}\) such that, whenever \(n\ge N_{j,k}\),
\[
\mathbb H_n^j\{|X_{j,n}-d_j|\ge k^{-1}\}<k^{-1}.
\]
Let \(\ell_j(n)\) be the largest \(k\le n\) with \(N_{j,k}\le n\), using \(\ell_j(n)=1\) when the set is empty, and set
\[
e_{j,n}=\ell_j(n)^{-1},
\qquad
\Gamma_n^j=\{\omega: |X_{j,n}(\omega)-d_j|<e_{j,n}\}.
\]
Then \(\ell_j(n)\to\infty\), so \(e_{j,n}\to0\), and
\[
\mathbb H_n^j(\Gamma_n^j)\to1,
\qquad
\sup_{\omega\in\Gamma_n^j}|G_n\sigma_n^2(Y_n^{j,\omega})-d_j|\to0 .
\]
These are the supports \(\Gamma_n^{\mathrm{same}}\) and \(\Gamma_n^{\mathrm{ind}}\). % lean: FiniteDesign.TendstoInProb.exists_uniform_support

\item Every full diagonal selection from either support belongs to \(\mathcal C_{\mathrm{dense}}\). Indeed, for the same-sign prior a selected row has the additive form
\[
Y_{i,n}^{\mathrm{same},\xi}(z,A)=\xi_i,\qquad \xi_i\in\{-1,1\},
\]
and for the independent-arm prior it has the additive form
\[
Y_{i,n}^{\mathrm{ind},\zeta}(z,A)=\zeta_{i,z},\qquad \zeta_{i,z}\in\{-1,1\}.
\]
Thus \(|Y_{i,n}(z,A)|=1\le B\). The original design skeleton supplies the group-count, treated-fraction, and sampling-fraction conditions in \cref{obj:def:dense-schedule-class}. For any diagonal selection with \(\omega_n\in\Gamma_n^j\) for every \(n\),
\[
G_n\sigma_n^2(Y_n^{j,\omega_n})\to d_j,
\]
because the support error is bounded by \(e_{j,n}\to0\). Since \(c_\sigma<d_j\), this convergence gives
\[
c_\sigma\le \liminf_{n\to\infty}G_n\sigma_n^2(Y_n^{j,\omega_n}).
\]
Hence all defining conditions of \(\mathcal C_{\mathrm{dense}}\) hold for each such diagonal selection. % lean: denseClass_withSchedule_of_scaled_tendsto

\item Since \(\mathbb H_n^j(\Gamma_n^j)\to1\) for \(j\in\{\mathrm{same},\mathrm{ind}\}\), both conditioning masses are positive for all sufficiently large \(n\). On those rows, let \(\mathbb E_{n,j}^{\Gamma}\) denote expectation under the one-realization mixture obtained by first conditioning \(\mathbb H_n^j\) on \(\Gamma_n^j\) and then applying the random partition and treatment design, and define
\[
\Delta_n=
\sup_{0\le \varphi\le1}
\left|
\mathbb E_{n,\mathrm{same}}^{\Gamma}\varphi(\mathcal O_n)
-
\mathbb E_{n,\mathrm{ind}}^{\Gamma}\varphi(\mathcal O_n)
\right|.
\]
For the finitely many remaining rows, take \(\Delta_n=0\); this totalization convention has no effect on the limiting assertion. For any bounded test \(0\le\varphi\le1\), conditioning a finite prior on a set \(\Gamma\) of positive prior probability changes its expectation by at most twice the complement probability:
\[
\left|
\mathbb E\{\varphi\mid \Gamma\}-\mathbb E\varphi
\right|
\le
2\{1-\mathbb H(\Gamma)\}.
\]
Indeed, writing \(m=\mathbb H(\Gamma)>0\), the conditioned expectation is
\[
\mathbb E\{\varphi\mid\Gamma\}=\frac{1}{m}\mathbb E\{\mathbf 1_\Gamma\varphi\},
\]
and hence
\[
\begin{aligned}
\left|
\mathbb E\{\varphi\mid \Gamma\}-\mathbb E\varphi
\right|
&=
\left|
\frac{1-m}{m}\mathbb E\{\mathbf 1_\Gamma\varphi\}
-
\mathbb E\{\mathbf 1_{\Gamma^c}\varphi\}
\right| \\
&\le
\frac{1-m}{m}\mathbb E\{\mathbf 1_\Gamma\varphi\}
+
\mathbb E\{\mathbf 1_{\Gamma^c}\varphi\} \\
&\le
\frac{1-m}{m}\,m+(1-m)
=2(1-m).
\end{aligned}
\]
This derivation also covers the case \(m=1\), where the unnormalized complement contribution is zero. Applying this bound to the two schedule priors on the eventual-positive rows and using the exact unconditioned equality from \cref{obj:prop:rademacher-mixture-separation} yields
\[
\Delta_n
\le
2\{1-\mathbb H_n^{\mathrm{same}}(\Gamma_n^{\mathrm{same}})\}
+
2\{1-\mathbb H_n^{\mathrm{ind}}(\Gamma_n^{\mathrm{ind}})\}.
\]
The right-hand side tends to \(0\), while \(\Delta_n\ge0\) because the zero test is admissible on the eventual-positive rows and by convention on the remaining rows. Therefore
\[
\Delta_n\to0.
\]
Together with the previous steps, this proves the support, dense-class, conditioned-total-variation, and uniform scaled-variance parts of the diagonal certificate. % lean: conditionedMixtureTV_le_complements

\item It remains to prove the all-statistic lower bound. Put
\[
\varepsilon=\frac{d_{\mathrm{same}}-d_{\mathrm{ind}}}{8d_{\mathrm{same}}},
\qquad
\eta=\frac{d_{\mathrm{same}}-d_{\mathrm{ind}}}{8}.
\]
Because \(0<c_\sigma<d_{\mathrm{ind}}<d_{\mathrm{same}}\), both \(\varepsilon\) and \(\eta\) are positive. Fix an arbitrary one-realization variance statistic \(\widehat S_n(\mathcal O_n)\), and define the threshold
\[
\vartheta=\frac{d_{\mathrm{same}}+d_{\mathrm{ind}}}{2}.
\]
For all sufficiently large \(n\), the supports are nonempty, have positive conditioning probability, and satisfy \(e_{\mathrm{same},n}<\eta\) and \(e_{\mathrm{ind},n}<\eta\). On such rows define the two wrong-decision tests
\[
\psi_n^{\mathrm{same}}(o)=\mathbf 1\{G_n\widehat S_n(o)\le\vartheta\},
\qquad
\psi_n^{\mathrm{ind}}(o)=\mathbf 1\{G_n\widehat S_n(o)>\vartheta\}.
\]
If \(\xi\in\Gamma_n^{\mathrm{same}}\) and
\[
\left|
\frac{\widehat S_n(o)}{\sigma_n^2(Y_n^{\mathrm{same},\xi})}-1
\right|\le\varepsilon,
\]
then the support bound gives
\[
|G_n\sigma_n^2(Y_n^{\mathrm{same},\xi})-d_{\mathrm{same}}|<\eta,
\]
and hence
\[
G_n\widehat S_n(o)
>
(1-\varepsilon)(d_{\mathrm{same}}-\eta)
>
\frac{d_{\mathrm{same}}+d_{\mathrm{ind}}}{2}.
\]
Thus \(\psi_n^{\mathrm{same}}(o)=1\) implies relative error greater than \(\varepsilon\). Similarly, if \(\zeta\in\Gamma_n^{\mathrm{ind}}\) and the relative error is at most \(\varepsilon\), then
\[
G_n\widehat S_n(o)
<
(1+\varepsilon)(d_{\mathrm{ind}}+\eta)
<
\frac{d_{\mathrm{same}}+d_{\mathrm{ind}}}{2},
\]
so \(\psi_n^{\mathrm{ind}}(o)=1\) implies relative error greater than \(\varepsilon\). The two strict algebraic margins are exactly the displayed choices of \(\varepsilon\) and \(\eta\). % lean: ratio_good_implies_same_decision

\item Let
\[
\mathfrak W_n(\widehat S,\varepsilon)
=
\sup_{Y\in\mathcal C_{\mathrm{dense}}}
\Pr_Y\!\left(
\left|
\frac{\widehat S_n(\mathcal O_n)}{\sigma_n^2(Y)}-1
\right|>\varepsilon
\right).
\]
For each \(\xi\in\Gamma_n^{\mathrm{same}}\), the preceding implication bounds the same-sign row probability of \(\psi_n^{\mathrm{same}}=1\) by the corresponding relative-error probability. That row can be embedded into a full triangular array in \(\mathcal C_{\mathrm{dense}}\): keep the row \(Y_n^{\mathrm{same},\xi}\) at index \(n\), choose arbitrary support rows \(\Gamma_m^{\mathrm{same}}\) for all sufficiently large \(m\ne n\), and choose arbitrary bounded sign rows at the finitely many remaining indices. The convergence and class argument from Step 3 applies because membership in the supports is eventual. Therefore
\[
\mathbb E_{n,\mathrm{same}}^{\Gamma}\psi_n^{\mathrm{same}}(\mathcal O_n)
\le
\mathfrak W_n(\widehat S,\varepsilon).
\]
The identical construction for \(\zeta\in\Gamma_n^{\mathrm{ind}}\) gives
\[
\mathbb E_{n,\mathrm{ind}}^{\Gamma}\psi_n^{\mathrm{ind}}(\mathcal O_n)
\le
\mathfrak W_n(\widehat S,\varepsilon).
\]
Averaging over the conditioned finite supports uses only positivity of the conditioning masses. % lean: tendsto_of_eventually_uniform_support_error

\item Since \(\psi_n^{\mathrm{same}}+\psi_n^{\mathrm{ind}}=1\) pointwise under the independent-arm observation channel,
\[
\mathbb E_{n,\mathrm{ind}}^{\Gamma}\psi_n^{\mathrm{same}}(\mathcal O_n)
=
1-\mathbb E_{n,\mathrm{ind}}^{\Gamma}\psi_n^{\mathrm{ind}}(\mathcal O_n).
\]
The definition of \(\Delta_n\), applied to \(\psi_n^{\mathrm{same}}\), gives
\[
\left|
\mathbb E_{n,\mathrm{same}}^{\Gamma}\psi_n^{\mathrm{same}}(\mathcal O_n)
-
\mathbb E_{n,\mathrm{ind}}^{\Gamma}\psi_n^{\mathrm{same}}(\mathcal O_n)
\right|
\le \Delta_n.
\]
Combining the last two displays yields
\[
1-\Delta_n
\le
\mathbb E_{n,\mathrm{same}}^{\Gamma}\psi_n^{\mathrm{same}}(\mathcal O_n)
+
\mathbb E_{n,\mathrm{ind}}^{\Gamma}\psi_n^{\mathrm{ind}}(\mathcal O_n).
\]
Using Step 6 on both terms,
\[
\frac{1-\Delta_n}{2}
\le
\mathfrak W_n(\widehat S,\varepsilon)
\]
for all sufficiently large \(n\). % lean: conditionedFiniteExpectation_compl

\item Finally, \(\Delta_n\to0\), so the left-hand side in the last display converges to \(1/2\). Also \(\mathfrak W_n(\widehat S,\varepsilon)\le1\): if the set of dense-class risks is nonempty, this is the supremum of probabilities bounded by one, and if it is empty the supremum convention used here gives \(0\). Hence
\[
\frac12
\le
\liminf_{n\to\infty}\mathfrak W_n(\widehat S,\varepsilon).
\]
Since \(\widehat S\) was arbitrary, the asserted lower bound holds for every measurable one-realization variance statistic. % lean: qv_diagonal_impossibility
\end{enumerate}
\end{proof}
